\section{RQ10: Does Syntactic Validity Guarantee Semantic Correctness?}
\label{appendix:rq10}

A core motivation for our execution-based evaluation is the observation that syntactic validity (compilation) is a necessary but insufficient condition for specification quality. In our experiments, we encountered cases where LLMs generated specifications that successfully compiled in Rocq but contained fundamental semantic flaws. These flaws often manifest as logical contradictions or overly restrictive constraints that rule out valid program behaviors.

This phenomenon highlights a critical failure mode: the model knows how to write valid Rocq code structure but fails to translate the natural language intent into a consistent logical reality. Because these specifications are syntactically perfect, they pass initial static checks, making them particularly deceptive. They only reveal their defects when instantiated against concrete test cases during the verification phase.

Figure \ref{fig:semantic_flaw_case} presents a representative case study from HumanEval/12 (\texttt{longest}). The task requires returning the first string of maximum length. The specification generated by Claude 4.5 Opus compiles without errors. However, it introduces a subtle logic error in how it defines the first occurrence. The specification inadvertently enforces a constraint that requires the result's index to be smaller than or equal to \textit{any} index of a max-length string. As shown in the trace, this logic collapses when the input list contains duplicate max-length strings (e.g., \texttt{["a", "a"]}), creating a mathematical contradiction ($1 \le 0$). This case underscores that specification generation is not just a translation task but requires rigorous semantic consistency, which LLMs still struggle to guarantee.

\begin{figure}[t!] % 增加放置参数，防止强行放在某处导致空间不足
    \centering
    % === 外层大盒子：蓝色背景 ===
    \begin{tcolorbox}[
        enhanced,
        title=\textbf{Case Study: Syntactically Valid but Semantically Incorrect Spec},
        fonttitle=\bfseries\small,
        colback=blue!3!white,    % 极淡蓝背景
        colframe=blue!50!black,  % 深蓝边框
        boxrule=0.8pt,
        arc=2pt,
        before skip=10pt,
        after skip=10pt
    ]
    
    % --- 1. Generated Specification 部分 ---
    \textbf{Generated Specification (Compiles Successfully):}
    \vspace{2pt}
    
    % 白色内盒子：无边框，纯白背景
    \begin{tcolorbox}[
        colback=white,
        colframe=white,
        frame hidden,       % 强制隐藏边框
        boxrule=0pt,
        sharp corners,
        top=2pt, bottom=2pt, left=2pt, right=2pt
    ]
    \begin{lstlisting}[
        language=Coq, 
        frame=none, 
        basicstyle=\ttfamily\fontsize{7.5pt}{8.5pt}\selectfont, % 微调字号确保显示全
        aboveskip=0pt,
        belowskip=0pt,
        breaklines=true,
        backgroundcolor={} 
    ]
Definition longest_spec (strings : list string) (result : option string) : Prop :=
  match strings with
  | [] => result = None
  | _ => let maxlen := max_length strings in
    exists s, result = Some s /\ In s strings /\ string_length s = maxlen /\
      (* Logic Error: Universal quantifier creates contradiction for duplicates *)
      (forall s'' idx1 idx2, 
        In s'' strings -> string_length s'' = maxlen ->
        nth_error strings idx1 = Some s ->   (* s is the result *)
        nth_error strings idx2 = Some s'' -> (* s'' is any max-len string *)
        idx1 <= idx2)
  end.
    \end{lstlisting}
    \end{tcolorbox}

    \vspace{0.2cm}
    
    % --- 2. Trace Analysis 部分 ---
    \textbf{Why the Logic Fails (Trace):}
    \vspace{2pt}
    
    \begin{tcolorbox}[
        colback=white,
        colframe=white,
        frame hidden,
        boxrule=0pt,
        sharp corners,
        top=2pt, bottom=2pt, left=2pt, right=2pt
    ]
    \begin{lstlisting}[
        language=Coq, 
        frame=none, 
        basicstyle=\ttfamily\fontsize{7.5pt}{8.5pt}\selectfont,
        aboveskip=0pt,
        belowskip=0pt,
        breaklines=true,
        escapechar=!,
        backgroundcolor={} 
    ]
(* Input: strings = ["a"; "a"], maxlen = 1. Implementation returns index 0 ("a"). *)
(* Spec requires condition for ALL index pairs pointing to max-len strings. *)
(* Note: "a" appears at both index 0 and index 1. *)

Variable idx1 := 1. (* "a" at index 1 *)
Variable idx2 := 0. (* "a" at index 0 *)

(* The Spec Constraint: idx1 <= idx2 *)
(* Substitution:        1 <= 0       *)
(* Result: False. Verification fails regardless of prover power. *)
    \end{lstlisting}
    \end{tcolorbox}

    \vspace{0.2cm}
    
    \end{tcolorbox}
    
    \caption{\textbf{Example of Semantic Error in HumanEval/12.} The model generates a Rocq specification that is syntactically correct but semantically flawed. It attempts to formalize the first longest string requirement but introduces a logic error that forbids duplicates. This results in a specification that is logically unsatisfiable for certain valid inputs, illustrating the gap between syntax generation and semantic understanding.}
    \label{fig:semantic_flaw_case}
\end{figure}